\documentclass[12pt, a4paper]{article}

\usepackage[margin=1.8cm,top=2cm,bottom=2cm]{geometry}
\usepackage{amsmath,amssymb,amsfonts}
\usepackage{graphicx}
\usepackage[colorlinks=true,linkcolor=blue,citecolor=blue,urlcolor=blue]{hyperref}
\usepackage{bm,xcolor}
\usepackage[numbers,sort&compress]{natbib}
\usepackage{float}

\newcommand{\Geff}{G_{\mathrm{eff}}}
\newcommand{\cs}{c_s}
\newcommand{\meff}{m_{\mathrm{eff}}}
\newcommand{\heff}{\hbar_{\mathrm{eff}}}
\newcommand{\rhob}{\bar{\rho}}

\begin{document}

\title{A Testable Prediction of the Logarithmic Superfluid Vacuum: Environment-Dependent Speed of Light}

\author{B. B. Kulangiev\\\small{Haifa, Israel}}
\date{April 2026}
\maketitle

\begin{abstract}
\noindent
In the logarithmic superfluid vacuum framework, the
speed of light is the speed of sound in the vacuum
condensate: $c_s^2 = [\ln(\rhob/\rho_c) + 2]^{-1}$.
Since the vacuum density $\rhob$ varies across
gravitational environments, the speed of light is not
a universal constant but an environment-dependent
quantity. We derive the variation
$\delta c/c = -GM/(2c^2 r)$ in the weak-field limit
and show that it is consistent with all existing
constraints, including the Cassini measurement
($|\gamma - 1| < 2.3 \times 10^{-5}$) and GW170817
($|c_{\mathrm{gw}} - c_{\mathrm{em}}|/c < 10^{-15}$).
The latter constrains only that gravitational and
electromagnetic waves propagate at the same speed
along the same path---which the framework predicts
exactly, since both are excitations of the same
condensate. The new, testable prediction is that
photons traversing cosmic voids (low vacuum density,
higher $c_s$) arrive systematically earlier than
photons traversing filaments (high vacuum density,
lower $c_s$). For a 60~Mpc void crossing, the
differential arrival time is of order seconds.
We identify three observational tests using existing
or near-future data: (i)~fast radio burst timing
residuals correlated with large-scale structure along
the sightline, (ii)~an additional contribution to CMB
temperature anisotropy beyond the integrated
Sachs-Wolfe effect, and (iii)~systematic differences
in time-of-flight for multi-messenger events from
different cosmic environments. A confirmed detection
would simultaneously validate the superfluid vacuum
framework and falsify the assumption of a strictly
constant speed of light.
\end{abstract}


% ====================================================================
\section{Introduction}
\label{sec:intro}
% ====================================================================

The constancy of the speed of light is a foundational
assumption of both special and general relativity. In
GR, the coordinate speed of light varies in a
gravitational field (the Shapiro delay), but the
locally measured speed is always $c$. No physical
mechanism exists within GR for the speed of light
to depend on the large-scale environment.

In the logarithmic superfluid vacuum
framework~\cite{kulangiev_2026a,kulangiev_2026b},
the speed of light is identified with the speed of
sound in the vacuum condensate. The sound speed
depends on the local condensate density through the
equation of state:
\begin{equation}
    c_s^2 = \frac{c^2}{\ln(\rhob/\rho_c) + 2}\,,
\label{eq:cs}
\end{equation}
where $\rhob$ is the local vacuum density and
$\rho_c = e\,\rho_\infty$ is a normalization constant
ensuring $c_s = c$ at the cosmological background
density $\rho_\infty$. Since $\rhob$ varies across
gravitational environments---higher near massive
objects, lower in cosmic voids---the speed of light
is an environment-dependent quantity.

This variation is not in conflict with special
relativity. Local observers measure $c_s$ using rulers
and clocks that are themselves made of condensate
excitations, and therefore co-vary with $c_s$. The
locally measured speed is always $c$, just as in
GR~\cite{will_2014}. The variation is detectable only
through \emph{differential} measurements: comparing
photon travel times across different environments.

In this paper, we derive the magnitude of the variation,
show that it is consistent with all existing constraints,
and identify three observational tests that can detect
or falsify the prediction using existing data.


% ====================================================================
\section{Derivation of the Speed Variation}
\label{sec:derivation}
% ====================================================================

\subsection{The weak-field expansion}

Writing $\rhob = \rho_\infty(1 + \epsilon)$ where
$\epsilon = \delta\rho/\rho_\infty$ is the vacuum
density perturbation, Eq.~(\ref{eq:cs}) becomes:
\begin{equation}
    c_s^2 = \frac{c^2}{1 + \ln(1+\epsilon)}\,.
\label{eq:cs_epsilon}
\end{equation}
For $|\epsilon| \ll 1$:
\begin{equation}
    c_s \approx c\left(1 - \frac{\epsilon}{2}\right)\,.
\label{eq:cs_linear}
\end{equation}
The fractional speed variation is:
\begin{equation}
    \frac{\delta c}{c} = -\frac{\epsilon}{2}
    = -\frac{1}{2}\frac{\delta\rho}{\rho_\infty}\,.
\label{eq:dc_c}
\end{equation}

\subsection{Connection to the gravitational potential}

From the Painlev\'{e}-Gullstrand flow
analysis~\cite{kulangiev_2026a,kulangiev_2026b}, the
vacuum density perturbation near a gravitating mass is
related to the dimensionless gravitational potential:
\begin{equation}
    \epsilon = \frac{v^2}{2c^2}
    = \frac{GM}{c^2 r} \equiv U\,,
\label{eq:epsilon_U}
\end{equation}
where we used the PG velocity
$v = \sqrt{2GM/r}$ and the Bernoulli
relation~\cite{kulangiev_2026b}. Therefore:
\begin{equation}
    \boxed{\;\frac{\delta c}{c}
    = -\frac{U}{2}
    = -\frac{GM}{2c^2 r}\;}
\label{eq:dc_U}
\end{equation}
The speed of light is lower near gravitating masses
(where the vacuum is denser) and higher in underdense
regions (cosmic voids). This is a monotonic function
of the gravitational potential.

\subsection{Numerical values across scales}

\begin{center}
\begin{tabular}{lcc}
\hline
Environment & $U = GM/(c^2 r)$ & $|\delta c/c|$ \\
\hline
Solar limb & $2.1 \times 10^{-6}$ & $1.1 \times 10^{-6}$ \\
Earth orbit & $1.0 \times 10^{-8}$ & $5.0 \times 10^{-9}$ \\
Solar neighbourhood (8 kpc) & $\sim 10^{-6}$ & $\sim 5 \times 10^{-7}$ \\
Cosmic void (30 Mpc, $\delta = -0.8$) & $\sim 10^{-5}$ & $\sim 5 \times 10^{-6}$ \\
KBC void (150 Mpc, $\delta = -0.5$) & $\sim 10^{-5}$ & $\sim 5 \times 10^{-6}$ \\
\hline
\end{tabular}
\end{center}

The variation is below $10^{-5}$ at all scales---well
within current measurement precision---but
accumulates over cosmological distances.


% ====================================================================
\section{Consistency with Existing Constraints}
\label{sec:constraints}
% ====================================================================

\subsection{The Cassini constraint}

The Cassini measurement~\cite{bertotti_2003}
constrains the PPN parameter $\gamma = 1.000021 \pm
0.000023$. In the superfluid framework, $\gamma = 1$
exactly for the PG flow~\cite{kulangiev_2026a}. The
speed variation $\delta c/c \sim 10^{-6}$ at the
solar limb does not affect $\gamma$ because $\gamma$
is determined by the ratio of spatial to temporal
metric perturbations, both of which scale with the
same $U$. The Cassini measurement constrains the
\emph{ratio}, not the individual variations.

\subsection{The GW170817 constraint}

The simultaneous detection of gravitational waves
and a gamma-ray burst from
GW170817~\cite{abbott_2017} constrains
$|c_{\mathrm{gw}} - c_{\mathrm{em}}|/c <
10^{-15}$. This is often cited as evidence that
the speed of light is exactly $c$.

However, this constraint applies only to the
\emph{same path through the same environment}. In
the superfluid framework, gravitational waves and
electromagnetic waves are both excitations of the
same condensate---phonons propagating on the same
acoustic metric. They travel at exactly the same
speed $c_s(\mathbf{x})$ along any given path:
\begin{equation}
    c_{\mathrm{gw}}(\mathbf{x})
    = c_{\mathrm{em}}(\mathbf{x})
    = c_s(\mathbf{x})\,.
\label{eq:cgw_cem}
\end{equation}
The GW170817 constraint is automatically
satisfied---it is a \emph{confirmation} of the
framework, not a constraint on it. The testable
prediction is different: $c_s$ varies across
environments, not between signal types.

\subsection{Local measurement invariance}

A local observer measures $c$ using apparatus made
of condensate excitations. If the local $c_s$ changes,
the ruler (atomic spacing $\propto \hbar/(\meff c_s)$)
and clock (atomic transition frequency
$\propto \meff c_s^2/\hbar$) change in concert, so
the measured ratio distance/time is always $c_s$
expressed in local units as a dimensionless number.
The variation is detectable only through comparisons
across different environments---precisely the
differential measurements described in the next
section.


% ====================================================================
\section{Observational Tests}
\label{sec:tests}
% ====================================================================

\subsection{Test 1: Fast radio burst timing residuals}

Fast radio bursts (FRBs) are millisecond-duration
radio transients at cosmological
distances~\cite{lorimer_2007,chatterjee_2017}. Their
dispersion measure (DM) provides the integrated
electron column density along the sightline. The
arrival time depends on both the DM (frequency-dependent
dispersion) and the integrated speed of light along
the path:
\begin{equation}
    t_{\mathrm{arr}} = \int_0^L \frac{dl}{c_s(l)}
    + t_{\mathrm{DM}}(\nu)\,.
\label{eq:t_arr}
\end{equation}
In GR, $c_s = c$ everywhere, and the first term is
simply $L/c$. In the superfluid framework:
\begin{equation}
    t_{\mathrm{arr}} = \frac{L}{c}
    + \frac{1}{c}\int_0^L \frac{U(l)}{2}\,dl
    + t_{\mathrm{DM}}(\nu)\,,
\label{eq:t_arr_svt}
\end{equation}
where the second term is the additional delay (or
advance) from the variable speed of light.

The prediction: FRBs whose sightlines pass through
cosmic voids should have negative timing residuals
(arrive earlier than DM-predicted), while those
passing through filaments or clusters should have
positive residuals (arrive later). The magnitude is:
\begin{equation}
    \delta t \sim \frac{L_{\mathrm{void}}}{c}
    \cdot \frac{|U_{\mathrm{void}}|}{2}
    \sim \text{seconds}
\label{eq:dt_frb}
\end{equation}
for a 60~Mpc void crossing.

This can be tested by cross-correlating the CHIME/FRB
catalog~\cite{chime_2021} with galaxy survey data
(SDSS, DESI) to classify sightlines as
void-dominated or filament-dominated, then comparing
the timing residuals of the two populations.

\subsection{Test 2: CMB anisotropy from variable \texorpdfstring{$c$}{c}}

CMB photons traversed the entire cosmic web between
last scattering and the present. Different photon
paths crossed different void and filament structures.
In GR, this produces the integrated Sachs-Wolfe (ISW)
effect---a temperature shift from the time-dependent
gravitational potential.

The superfluid framework predicts an
\emph{additional} contribution: photons that
traversed voids traveled faster (higher $c_s$) and
therefore arrive with a slightly different frequency
than those that traversed filaments. This creates a
temperature anisotropy:
\begin{equation}
    \frac{\delta T}{T} \sim
    \frac{\delta c}{c} \sim \frac{U}{2}
    \sim 10^{-5}\,,
\label{eq:dT_T}
\end{equation}
correlated with the large-scale structure. This is
of the same order as the ISW effect and may be
partially degenerate with it, but the angular
dependence differs: the variable-$c$ contribution
depends on the \emph{integral} of the potential
along the path, while the ISW effect depends on its
\emph{time derivative}.

A targeted stacking analysis---correlating CMB
temperature with known void and supercluster
catalogs~\cite{sutter_2012}---could separate the two
contributions.

\subsection{Test 3: Multi-messenger events from
different environments}

If two neutron star mergers occur at the same
redshift but in different large-scale environments
(one in a void, one in a cluster), the GW and EM
signals from each event should arrive with the same
$c_{\mathrm{gw}} = c_{\mathrm{em}}$ (satisfying
GW170817-type constraints). However, the
\emph{absolute} arrival times of the two events
(referenced to their cosmological distances) should
differ by:
\begin{equation}
    \Delta t \sim \frac{D}{c}
    \cdot \frac{|U_{\mathrm{void}}
    - U_{\mathrm{cluster}}|}{2}\,,
\label{eq:dt_mm}
\end{equation}
where $D$ is the luminosity distance. For events at
$D \sim 100$~Mpc with $\Delta U \sim 10^{-5}$, this
gives $\Delta t \sim 0.1$~s---detectable with
LIGO/Virgo timing precision.

This test requires multiple multi-messenger events
at similar redshifts but different environments,
which the upcoming LIGO O5 run and Vera Rubin
Observatory may provide.


% ====================================================================
\section{The Nonlinear Regime: Cosmological Implications}
\label{sec:nonlinear}
% ====================================================================

The linear result $\delta c/c = -U/2$ applies in the
weak-field limit $|\epsilon| \ll 1$. At cosmological
scales, the full nonlinear expression~(\ref{eq:cs_epsilon})
becomes relevant.

For a deep void with $\epsilon = -0.8$ (80\%
underdense):
\begin{equation}
    c_s = \frac{c}{\sqrt{1 + \ln(0.2)}}
    = \frac{c}{\sqrt{1 - 1.609}}
    = \frac{c}{\sqrt{-0.609}}\,,
\end{equation}
which is imaginary---indicating that the linear
equation of state breaks down for $\epsilon < 1 - e
\approx -0.632$. This is the \emph{vacuum cavitation
threshold}: below this density, the condensate cannot
support sound waves, and the acoustic metric
description fails.

This has two implications:

(i) Cosmic voids cannot be arbitrarily deep. The
maximum void depth is $\delta\rho/\rho_0 < 1 - 1/e
\approx 0.632$, set by the equation of state rather
than by gravitational dynamics alone. Observed voids
with $\delta \sim -0.8$ to $-0.9$ would violate this
bound, suggesting either that the linear
approximation overestimates the baryonic density
contrast, or that the vacuum density perturbation is
smaller than the matter density perturbation (as
expected, since the vacuum is much stiffer than
matter).

(ii) Near the cavitation threshold, $c_s$ increases
rapidly---photons traversing the deepest voids
experience a significant speedup. This nonlinear
amplification could enhance the differential timing
signal beyond the linear estimate, potentially making
it detectable with current instruments.


% ====================================================================
\section{Discussion}
\label{sec:discussion}
% ====================================================================

\subsection{What is predicted and what is not}

The framework predicts:
\begin{itemize}
\item $c_{\mathrm{gw}} = c_{\mathrm{em}}$ exactly
  (both are condensate excitations). Confirmed by
  GW170817.
\item $c_s$ varies across gravitational environments
  by $\delta c/c \sim U/2$. Not yet tested.
\item Photon arrival times correlate with large-scale
  structure. Testable with FRBs + galaxy surveys.
\end{itemize}

The framework does \emph{not} predict:
\begin{itemize}
\item A frequency-dependent speed of light. In the
  acoustic limit ($\omega \ll \meff c^2/\hbar$), the
  dispersion relation is linear:
  $\omega = c_s |\mathbf{k}|$. Frequency dependence
  appears only near the Planck scale
  ($\omega \sim \meff c^2/\hbar$), where the Bohm
  potential modifies the dispersion
  relation~\cite{amelino_2013}.
\item A violation of Lorentz invariance. The acoustic
  metric is locally Lorentzian with speed $c_s$. Local
  physics is Lorentz-invariant; only the
  \emph{cosmological} background breaks the symmetry
  by providing a preferred frame (the condensate rest
  frame $\equiv$ the CMB frame).
\end{itemize}

\subsection{Relation to other variable-\texorpdfstring{$c$}{c} proposals}

Variable speed of light (VSL) theories have been
proposed by Moffat~\cite{moffat_1993} and
Magueijo~\cite{magueijo_2003} as alternatives to
inflation. These typically modify the Maxwell or
Einstein action to allow $c$ to vary as a function
of cosmic time.

The present proposal differs in three respects:
(i)~$c$ varies with \emph{spatial environment}, not
cosmic time; (ii)~the variation is derived from a
specific equation of state, not postulated;
(iii)~both $c_{\mathrm{gw}}$ and
$c_{\mathrm{em}}$ vary together, preserving
$c_{\mathrm{gw}} = c_{\mathrm{em}}$.

\subsection{Spatial variation of the gravitational
constant}

Because the effective gravitational constant in
this framework depends directly on the sound speed
($\Geff = \eta c_s^2/(4\pi\rho_0)$), the
environment-dependent variation of $c_s$ implies
that $G$ must also vary across cosmological
environments. While the local PG flow shields
individual solar systems from this variation
($\epsilon_1 = 0$), the global difference in
$\rho_0$ between cosmic voids and filaments
leads to a macroscopic spatial variation in $G$.
The full cosmological implications of this
variation---including the quantitative resolution
of the Hubble tension via the
$\Lambda \propto c^4$ lever arm and a
hydrodynamic realization of Mach's
principle---are derived in the companion
paper~\cite{kulangiev_2026_cosmo}.

\subsection{Falsifiability}

The prediction is cleanly falsifiable. If FRB timing
residuals show \emph{no} correlation with large-scale
structure after controlling for DM and scattering,
the framework's prediction of environment-dependent
$c_s$ is falsified. This is a binary test: the
correlation either exists or it does not. No
parameter adjustment can remove the prediction
while preserving the framework.


% ====================================================================
\section{Conclusion}
\label{sec:conclusion}
% ====================================================================

The logarithmic superfluid vacuum framework makes a
specific, testable prediction that distinguishes it
from general relativity: the speed of light varies
across gravitational environments as
$\delta c/c = -GM/(2c^2 r)$. This variation is
consistent with all existing constraints (Cassini,
GW170817, local Lorentz invariance tests) because
these measurements either probe the same environment
or compare signals that co-propagate on the same
acoustic metric.

The new, falsifiable prediction is that photons
traversing cosmic voids arrive systematically earlier
than those traversing filaments, by an amount
proportional to the integrated gravitational potential
difference along the path. Three observational tests
are identified:

(1) FRB timing residuals correlated with large-scale
structure (testable with CHIME/FRB + SDSS/DESI).

(2) CMB temperature anisotropy from variable $c_s$
beyond the ISW effect (testable with Planck data +
void catalogs).

(3) Systematic arrival time differences for
multi-messenger events from different environments
(testable with LIGO O5 + Vera Rubin).

A confirmed detection would validate the superfluid
vacuum framework and demonstrate that the speed of
light is not a universal constant but the sound speed
of the medium in which we are immersed.


% ====================================================================
\section*{Acknowledgments}
% ====================================================================

The author thanks K.~G.~Zloshchastiev for foundational
work on the logarithmic superfluid vacuum.
 Claude (Anthropic, claude-opus-4-6) and Gemini (Google, Gemini 3.1 Pro) were used as editorial and computational tools during manuscript preparation. All scientific content, equations, and
interpretations are the author's own.


\begin{thebibliography}{20}

\bibitem{kulangiev_2026a}
B.~B.~Kulangiev,
Conditions for Emergent Gravitational Light Bending from a Logarithmic Superfluid Vacuum,
Preprint (2026).
\url{https://kulangiev.github.io#paper1}

\bibitem{kulangiev_2026b}
B.~B.~Kulangiev,
Emergent Newtonian Gravity and the Gravitational Feynman-Onsager Relation in the Logarithmic Superfluid Vacuum,
Preprint (2026).
\url{https://kulangiev.github.io#combined}

\bibitem{kulangiev_2026_cosmo}
B.~B.~Kulangiev,
Cosmological Implications of the Logarithmic Superfluid Vacuum: Dark Energy, the Hubble Tension, and a Vacuum Self-Consistency Equation,
Preprint (2026).
\url{https://kulangiev.github.io#paper3}

\bibitem{will_2014}
C.~M.~Will,
\emph{Theory and Experiment in Gravitational Physics},
2nd ed.\ (Cambridge, 2014).

\bibitem{bertotti_2003}
B.~Bertotti, L.~Iess, and P.~Tortora,
A test of general relativity using radio links with
the Cassini spacecraft,
Nature \textbf{425}, 374 (2003).

\bibitem{abbott_2017}
B.~P.~Abbott \emph{et al.},
Gravitational waves and gamma-rays from a binary
neutron star merger: GW170817 and GRB~170817A,
Astrophys.\ J.\ Lett.\ \textbf{848}, L13 (2017).

\bibitem{lorimer_2007}
D.~R.~Lorimer, M.~Bailes, M.~A.~McLaughlin,
D.~J.~Narkevic, and F.~Crawford,
A bright millisecond radio burst of extragalactic
origin,
Science \textbf{318}, 777 (2007).

\bibitem{chatterjee_2017}
S.~Chatterjee \emph{et al.},
A direct localization of a fast radio burst and its
host,
Nature \textbf{541}, 58 (2017).

\bibitem{chime_2021}
CHIME/FRB Collaboration,
The first CHIME/FRB fast radio burst catalog,
Astrophys.\ J.\ Suppl.\ \textbf{257}, 59 (2021).

\bibitem{sutter_2012}
P.~M.~Sutter, G.~Lavaux, B.~D.~Wandelt, and
D.~H.~Weinberg,
A public void catalog from the SDSS DR7 galaxy
redshift surveys,
Astrophys.\ J.\ \textbf{761}, 44 (2012).

\bibitem{amelino_2013}
G.~Amelino-Camelia,
Quantum-spacetime phenomenology,
Living Rev.\ Rel.\ \textbf{16}, 5 (2013).

\bibitem{moffat_1993}
J.~W.~Moffat,
Superluminary universe: A possible solution to the
initial value problem in cosmology,
Int.\ J.\ Mod.\ Phys.\ D \textbf{2}, 351 (1993).

\bibitem{magueijo_2003}
J.~Magueijo,
New varying speed of light theories,
Rep.\ Prog.\ Phys.\ \textbf{66}, 2025 (2003).

\end{thebibliography}

\end{document}
